Fine-tuning a code model on obfuscated programs raises its accuracy on those programs by roughly
nineteen points. We ask what that number is made of: competence at reading the program, or
competence at the surface a particular transform leaves behind. We separate the two with tests that
share no machinery --- whether the gain survives being stacked with a transform family the model has
never seen, and whether it survives the task being run backwards. The field's existing answers,
inference-time routing, weight-space merging, and training on a breadth of transforms, fail both.
Breadth composes on stacks built from what it saw and loses that advantage the moment an unseen
family enters, a gap significant on every one of six models across four lineages; and every method
that fits the forward task buys its gain by losing backward competence on the very same programs.
Three further instruments locate what is learned instead: the surface a transform leaves, which is
exactly what an obfuscator is free to change. That diagnosis names its own remedy. We anchor the
model to the one thing a semantics-preserving transform leaves invariant, its own behaviour on the
unobfuscated parent of each program, through a paired-consistency objective that trains on exactly
the rows breadth trains on. Re-running the same two tests, the anchored model is the only one whose
advantage \emph{grows} as input diverges from training, rising above the clean-code control on
stacks containing an unseen family, and the only one whose forward gain costs nothing backwards. Its
repair holds on six of eight models and reverses on one, and on the unseen family alone it is as
robust as clean-code tuning and no more. Adaptation to obfuscated code buys recombination of what
the data covered; anchoring recovers what that costs; neither buys invariance.
